% X_0^21(58)* is hyperelliptic.  The proof is the counting argument (CM divisors of degree two that
% are pairwise torsion-related, against the torsion bound), with the torsion relations certified by
% a dual test against weight-3/2 forms.  At a level with 2 || N the dual test needs a lemma relating
% the Heegner series of X_0^D(N)* to that of X_0^D(N/2)*, proved here.  Tests:
% tests/Hyperelliptic21_58.m and tests/_offline/Hyperelliptic21_58.m; solver in tools/proof-21-58/.
\documentclass[11pt,reqno]{amsart}
\usepackage[margin=1.15in]{geometry}
\usepackage{amsmath,amssymb,amsthm,mathtools,booktabs}
\usepackage[colorlinks=true,linkcolor=blue,citecolor=blue,urlcolor=blue]{hyperref}
\theoremstyle{plain}
\newtheorem{theorem}{Theorem}[section]
\newtheorem{proposition}[theorem]{Proposition}
\newtheorem{lemma}[theorem]{Lemma}
\newtheorem{corollary}[theorem]{Corollary}
\theoremstyle{definition}
\newtheorem{remark}[theorem]{Remark}
\DeclareMathOperator{\Jac}{Jac}
\DeclareMathOperator{\Pic}{Pic}
\DeclareMathOperator{\CM}{CM}
\DeclareMathOperator{\Div}{Div}
\DeclareMathOperator{\nm}{nm}
\DeclareMathOperator{\tr}{tr}
\newcommand{\Q}{\mathbb{Q}}
\newcommand{\Z}{\mathbb{Z}}
\newcommand{\F}{\mathbb{F}}
\newcommand{\HH}{\mathfrak{H}}
\newcommand{\cO}{\mathcal{O}}
\newcommand{\Hs}{\mathcal{H}}
\begin{document}
\title{$X_0^{21}(58)^*$ is hyperelliptic}
\author{}
\date{\today}
\maketitle

\section{Statement}

Let $X=X_0^{21}(58)$, the Shimura curve of an Eichler order of level $58$ in the quaternion algebra
of discriminant $21$, and $C=X^*=X/W$ its quotient by the full Atkin--Lehner group, of order $16$.
Then $C$ has genus three, its Jacobian is simple and isogenous to $A_f$ for the newform
$f=\texttt{609.2.a.d}$ of level $609=21\cdot29$ (so $f$ is old at $2$ on $X$), and no Atkin--Lehner
involution survives on $C$.

\begin{theorem}\label{thm:main}
$C=X_0^{21}(58)^*$ is hyperelliptic. The fibres of its hyperelliptic map include
\[
F_2=\{P_{-16},P_{-120}\},\quad F_4=\{P_{-88},P_{-100}\},\quad F_5=\{P_{-112},P_{-372}\},\quad
F_6=\{P_{-168},P_{-228}\},
\]
$P_d$ the rational CM point of discriminant $d$ on $C$.
\end{theorem}

The argument is the one of \texttt{paper/torsion-counting-hyperelliptic.tex}: the classes of the four
effective divisors $F_i$ of degree two lie in one coset of $\Jac(C)(\Q)_{\mathrm{tors}}$, which has
order at most $2$, so two of them coincide, and two distinct effective divisors of degree two in one
class give a $g^1_2$. What is new is how the torsion relations $F_i-F_j$ are certified. On the
curves treated so far they came from Borcherds forms given by their principal parts; at level $2436$
with $2\,\|\,N$ that formulation is blind to half of the relevant modular forms, and the certificate
goes instead through the curve $X_0^{21}(29)^*$ of half the level. The lemma that makes this work
(Section~\ref{sec:lemma}) is general: it applies to every $X_0^D(N)^*$ with $D$ odd and $N=2N'$,
$N'$ odd.

\section{The dual test}\label{sec:dual}

We recall the setting of \cite{gy} and of \texttt{paper/hyperelliptic-6-379.tex}. For an Eichler
order $\cO$ of level $N$ in the quaternion algebra of discriminant $D$ let $L=\{x\in\cO:\tr x=0\}$,
a lattice of level $M=4\cdot\mathrm{oddpart}(DN)$, and let $G=N_B^+(\cO)/\Q^\times$, so that
$X^*=G\backslash\HH$. A weakly holomorphic form $f=\sum_nc(n)q^n$ of weight $\tfrac12$ on
$\Gamma_0(M)$ with the theta multiplier, holomorphic at every cusp other than $\infty$, induces a
vector-valued form $F_f$ of type $\rho_L$ whose principal part is $\sum_{n>0}c(-n)q^{-n}e_0$
\cite[Lemmas 10 and 24]{gy}, and the Borcherds lift $\Psi_f$ of $F_f$ is a meromorphic modular form
on $X^*$ of weight $c(0)$ with divisor
\begin{equation}\label{eq:div}
\operatorname{div}\Psi_f=\sum_{n>0}c(-n)\,Z(n),\qquad
Z(n)=\sum_{\substack{r\ge1\\ -4n/r^2\ \text{a discriminant}}}\frac1{e_{-4n/r^2}}\,P_{-4n/r^2},
\end{equation}
where $P_d$ is the sum of the CM points of discriminant $d$ on $X^*$ (zero if there are none) and
$e_d$ the order of their stabiliser in $G$ \cite[Lemmas 22 and 25]{gy}.

\begin{proposition}\label{prop:dual}
Let $\lambda\in\Pic(X^*)_\Q$ be the Hodge class, and for $\varphi$ a linear form on $\Pic(X^*)_\Q$
with $\varphi(\lambda)=0$ put $\Phi_\varphi=\sum_{n\ge1}\varphi(Z(n))q^n$. Then $\Phi_\varphi$ is the
$q$-expansion of a cusp form of weight $\tfrac32$ on $\Gamma_0(M)$ with the multiplier of
$\theta^3$.
\end{proposition}

\begin{proof}
Let $c=(c(-n))_{n\ge1}$ be a finitely supported vector. There is a form $f$ as above with principal
part $\sum c(-n)q^{-n}$ if and only if $\sum_nc(-n)a_g(n)=0$ for every cusp form $g$ of weight
$\tfrac32$ on $\Gamma_0(M)$ with the conjugate multiplier: the differential $fg\,d\tau$ on $X_0(M)$
has its only pole at $\infty$, with residue $\sum c(-n)a_g(n)$, and Serre duality makes the
vanishing of these residues sufficient. For such $f$, \eqref{eq:div} gives $\sum c(-n)Z(n)=k\lambda$
in $\Pic(X^*)_\Q$ with $k$ the weight of $\Psi_f$, so $\sum c(-n)\varphi(Z(n))=0$. Thus
$\Phi_\varphi$ pairs to zero with every vector in the annihilator of $S_{3/2}(\Gamma_0(M))$ under
the residue pairing; by finite dimensionality it is itself a cusp form.
\end{proof}

The \emph{dual test} follows. A degree-zero divisor $\Delta$ supported on CM points of $X^*$ is
written as $\sum c(-n)Z(n)$ by solving, for every discriminant reachable from the poles considered,
the linear condition that the coefficients sum to $e_d$ times the prescribed multiplicity. If $V$ is a
space of formal $q$-series containing every $\Phi_\varphi$ and $c$ is orthogonal to $V$, then
$\varphi(\Delta)=0$ for all $\varphi$, i.e.\ $\Delta$ is torsion in $\Jac(X^*)(\Q)$. The natural
choice $V=W:=\{g\in S_{3/2}:a_g(n)=0$ whenever $Z(n)=0\}$ is sharp when $2\mid D$: there the
$2$-adic condition for CM points removes about half of all indices and $W$ has dimension equal to
the rank of $\Jac(X^*)(\Q)$. When $2\,\|\,N$ almost every index carries CM points and $W$ is twice
too large; the lemma below cuts it to size.

\section{The lemma for $2\,\|\,N$}\label{sec:lemma}

Throughout $D$ is odd, $N=2N'$ with $N'$ odd and squarefree, $\cO=\cO_N\subset\cO_{N'}$ are Eichler
orders of levels $N$ and $N'$ with $\cO_N\subset\cO_{N'}$, $G=N^+(\cO_N)/\Q^\times$ and
$G'=N^+(\cO_{N'})/\Q^\times$, $X^*=G\backslash\HH$, $X'^*=G'\backslash\HH$. Both lattices have level
$M=4DN'$. For an order $R$ of discriminant $d$ in an imaginary quadratic field put
\[
\chi_d(2)=\begin{cases}\left(\tfrac d2\right)&\text{if }2\nmid\text{cond}(R),\\ 0&\text{if }2\mid\text{cond}(R),\end{cases}
\]
and write $R_2=\Z+2R$, of discriminant $4d$.

\subsection{Divisors on the upper half plane}
For a subgroup $\Gamma\subset\mathrm{PGL}_2^+(\mathbb{R})$ acting on $\HH$ with finite stabilisers,
a $\Gamma$-invariant locally finite divisor $\sum_\tau m_\tau[\tau]$ on $\HH$ defines the
$\Q$-divisor $\sum_{x}(m_\tau/e_\tau)[x]$ on the coarse curve $\Gamma\backslash\HH$, $e_\tau$ the
order of the stabiliser of $\tau$; this is how \eqref{eq:div} arises, the divisor of $\Psi_f$ on
$\HH$ having multiplicity $\sum c(-n)$ at each CM point. For $\Gamma_1\subset\Gamma_2$ of finite
index and the coarse map $\pi$, pullback of $\Q$-divisors is the identity on invariant divisors of
$\HH$, and pushforward is $D\mapsto\sum_{g\in\Gamma_2/\Gamma_1}g\cdot D$ (one checks both against
the ramification indices $e^{\Gamma_2}_\tau/e^{\Gamma_1}_\tau$ of the coarse map).

For $\tau\in\HH$ let $K_\tau\subset B$ be the quadratic field fixing $\tau$ and $R^{(N)}_\tau=K_\tau\cap\cO_N$,
$R'_\tau=K_\tau\cap\cO_{N'}$. Put
\[
C(d)=\sum_{\tau:\ R^{(N)}_\tau\simeq R_d}[\tau],\qquad C'(d)=\sum_{\tau:\ R'_\tau\simeq R_d}[\tau],
\]
$G$- resp.\ $G'$-invariant divisors on $\HH$, which define $e_d^{-1}P_d$ on $X^*$ and
$e_d'^{-1}P'_d$ on $X'^*$. Thus $Z(n)=\sum_rC(-4n/r^2)$, and we set
\[
Z'(n)=\sum_{r}C'(-4n/r^2),\qquad H(n)=\sum_{\substack{s\ge1\\ -n/s^2\ \text{a discriminant}}}C'(-n/s^2).
\]
$H(n)$ is nonzero only for $n\equiv0,3\pmod4$, and $H(4n)=Z'(n)$, since $-4n/s^2$ runs over all
$s\ge1$. Also $Z'(n)=O'(n)+H(n)$ with $O'(n)=\sum_{r\text{ odd}}C'(-4n/r^2)$.

\subsection{The correspondence}
Let $G''=G\cap G'$. The elements of $G'$ normalising $\cO_N$ are $\Gamma_NW_{N'}$ (the normaliser
of an Eichler order of level $2$ in $\mathrm{GL}_2(\Z_2)$ is the order itself), so
$G''=\Gamma_NW_{N'}$, of index $2$ in $G$ (missing $w_2$) and of index $3$ in $G'$. Let
$Y=G''\backslash\HH$ with $q:Y\to X^*$ and $p:Y\to X'^*$, and $T=p_*q^*$ on $\Q$-divisors.

\begin{lemma}\label{lem:fibre}
$T\,C(d)=(1+\chi_d(2))\,C'(d)+(2-\chi_{d/4}(2))\,C'(d/4)$, the second term being zero if $d/4$ is
not a discriminant.
\end{lemma}

\begin{proof}
On divisors of $\HH$, $q^*$ is the identity and $p_*C(d)=\sum_{i=1}^3g_i\,C(d)$ for coset
representatives $g_i$ of $G'/G''$, which may be taken in $\Gamma_{N'}$. Now $g_iC(d)$ is the sum of
$[\tau]$ over $\tau$ with $K_\tau\cap g_i\cO_Ng_i^{-1}\simeq R_d$, and $g_i\cO_Ng_i^{-1}$, $i=1,2,3$,
are the three Eichler orders of level $N$ inside $\cO_{N'}$: they correspond to the three lines of
$\F_2^2$ through $\cO_{N'}/2\cO_{N'}\simeq M_2(\F_2)$, permuted transitively by $\Gamma_{N'}$. Fix
$\tau$ with $R'_\tau=R$ of discriminant $\tilde d$. The order $K_\tau\cap\cO^{(i)}$ is the
stabiliser in $R$ of the $i$-th line, which is $R$ if the line is $R$-stable and $\Z+2R=R_2$
otherwise: an element of $R$ fixing a non-stable line acts on $\F_2^2$ as a scalar, and conversely.
The number of $R$-stable lines is $1+\chi_{\tilde d}(2)$: $R/2R$ is $\F_2\times\F_2$, $\F_4$ or
$\F_2[\varepsilon]/(\varepsilon^2)$ according as $2$ splits, is inert, or is ramified in $R$, with
$2$, $0$, $1$ stable lines, and when $2\mid\text{cond}(R)$ the element $2\omega_0$ ($\omega_0$ a
generator of the maximal order) acts as a nonzero nilpotent, with one stable line. So $\tau$ occurs
in $p_*C(d)$ with multiplicity $1+\chi_d(2)$ if $R'_\tau\simeq R_d$, and with multiplicity
$2-\chi_{d/4}(2)$ if $R'_\tau\simeq R_{d/4}$, and not otherwise.
\end{proof}

\begin{lemma}\label{lem:series}
$T\,Z(n)=H(4n)+2H(n)$ for every $n\ge1$; in other words $T\circ Z=(U(4)+2)H$ as series of divisors
on $X'^*$.
\end{lemma}

\begin{proof}
Apply Lemma~\ref{lem:fibre} to each term $C(-4n/r^2)$ of $Z(n)$. For $r$ odd, $d=-4n/r^2\equiv0\pmod4$
has $\chi_d(2)=0$: either $d/4$ is a discriminant and then $2\mid\text{cond}(R_d)$, or
$d/4\equiv2,3\pmod4$ and then $2$ ramifies in $R_d$. For $r=2s$, $d=-n/s^2$. Hence
\[
T\,Z(n)=\sum_{r\text{ odd}}C'(-4n/r^2)+\sum_{s}\bigl(1+\chi_{-n/s^2}(2)\bigr)C'(-n/s^2)
+\sum_{s}\bigl(2-\chi_{-n/s^2}(2)\bigr)C'(-n/s^2),
\]
the last two sums over $s$ with $-n/s^2$ a discriminant (the third comes from the $d/4$ terms of
all $r$, odd and even, with $d/4=-n/r^2$). The characters cancel, leaving $O'(n)+3H(n)=Z'(n)+2H(n)=H(4n)+2H(n)$.
\end{proof}

\begin{lemma}\label{lem:inj}
Suppose $\Jac(X^*)$ is isogenous to $A_f$ for a newform $f$ of level $DN'$ (new at the primes of
$D$, old at $2$). Then $T:\Jac(X^*)_\Q\to\Jac(X'^*)_\Q$ is injective.
\end{lemma}

\begin{proof}
Let $\pi_1=p$ and $\pi_2=p\circ w_2:Y\to X'^*$, the two degeneracy maps. The $f$-part of $\Jac(Y)$ is
$\pi_1^*A_f\oplus\pi_2^*A_f$, and $q^*$ identifies $\Jac(X^*)_\Q$ with its $w_2$-invariants
$\{\pi_1^*y+\pi_2^*y\}$. Then $T(\pi_1^*y+\pi_2^*y)=3y+(\pi_1)_*\pi_2^*y=(3+T_2)y$ with $T_2$ the
Hecke correspondence at the good prime $2$ of $X'^*$, which acts on $A_f$ through the eigenvalue
$a_2(f)$ of absolute value at most $2\sqrt2<3$.
\end{proof}

\begin{proposition}\label{prop:plus}
For every linear form $\varphi'$ on $\Pic(X'^*)_\Q$ killing the Hodge class, $H_{\varphi'}:=\sum_n\varphi'(H(n))q^n$
is the $q$-expansion of a cusp form in Kohnen's plus space $S^+_{3/2}(\Gamma_0(M))$, i.e.\ a form
of weight $\tfrac32$ on $\Gamma_0(M)$ with the multiplier of $\theta^3$ whose coefficients vanish
unless $n\equiv0,3\pmod4$.
\end{proposition}

\begin{proof}
By \cite[Corollary 9]{gy}, $L'^\vee/L'\simeq\Z/2\times(\Z/DN')^2$; the $\Z/2$ is generated by
$\ell_1/2$ with $\nm(\ell_1/2)\equiv-\tfrac14\pmod1$, so the $2$-part of the discriminant form is
that of a rank-one lattice and $\rho_{L'}=\rho_2\otimes\rho_{\mathrm{odd}}$ accordingly. By
Borcherds \cite[Theorem 4.5]{borcherds}, the generating series
$\mathcal F=\sum_{\eta}\sum_{m}\varphi'(y_{m,\eta})q^me_\eta$ of the Heegner divisors of $L'$ is a
vector-valued modular form of weight $\tfrac32$ and type $\rho_{L'}$ (for the dual sign convention),
cuspidal because $\varphi'$ kills the Hodge class. Its component at $(0,0)$ has coefficients
$\varphi'(Z'(n))$, and its component at $(\ell_1/2,0)$ has coefficients $\varphi'(\sum_{r\text{ odd}}C'(-m/r^2))$
at exponent $m/4$ with $m\equiv3\pmod4$: a vector $\lambda\in L'+\ell_1/2$ has $\tfrac12+\lambda\in\cO_{N'}$
and $\lambda^\perp$ is a CM point of odd discriminant, as in the proof of \cite[Lemma 25]{gy}. For
$\gamma\in\Gamma_0(DN')$ the Weil representation of the odd part satisfies
$\rho_{\mathrm{odd}}(\gamma)e_0=\chi(\gamma)e_0$ for a character $\chi$ (see e.g.\ \cite[Theorem 4.7]{scheithauer}),
so the pair of components $(\mathcal F_{(0,0)},\mathcal F_{(\ell_1/2,0)})$ is a vector-valued form
of type $\rho_2\otimes\chi$ on $\Gamma_0(DN')$, and the map
$(F_0,F_{1/2})\mapsto F_0(4\tau)+F_{1/2}(4\tau)$ identifies such forms with the Kohnen plus space
of level $4DN'$ (\cite[\S5]{ez} for level one, \cite{kohnen} for odd level). The image is exactly $H_{\varphi'}$: its
coefficient at $n\equiv0\pmod4$ is $\varphi'(Z'(n/4))=\varphi'(H(n))$ and at $n\equiv3\pmod4$ it is
$\varphi'(\sum_{r\text{ odd}}C'(-n/r^2))=\varphi'(H(n))$. The character $\chi$ is trivial here,
because the scalar forms of \cite[Lemma 10]{gy} that induce $\rho_{L'}$-forms carry the bare theta
multiplier.
\end{proof}

\begin{theorem}[the test for $2\,\|\,N$]\label{thm:test}
Assume Lemma~\ref{lem:inj}. Let
\[
W^+=\{g\in S_{3/2}(\Gamma_0(M)):a_g(n)=0\ \text{unless }n\equiv0,3\ (4)\text{ and }H(n)\neq0\},\qquad
V=(U(4)+2)\,W^+ .
\]
If a degree-zero divisor $\Delta=\sum_nc(-n)Z(n)$ on $X^*$ has $\sum_nc(-n)v(n)=0$ for all $v\in V$,
then $\Delta$ is torsion in $\Jac(X^*)(\Q)$.
\end{theorem}

\begin{proof}
Every $\varphi$ on $\Jac(X^*)_\Q$ is $\varphi'\circ T$ for some $\varphi'$, by Lemma~\ref{lem:inj}.
By Lemma~\ref{lem:series}, $\varphi(Z(n))=\varphi'(H(4n))+2\varphi'(H(n))$, i.e.\
$\Phi_\varphi=(U(4)+2)H_{\varphi'}$, and $H_{\varphi'}\in W^+$ by Proposition~\ref{prop:plus} and
because $H_{\varphi'}(n)=0$ when $H(n)=0$. So $\varphi(\Delta)=\sum c(-n)\Phi_\varphi(n)=0$.
\end{proof}

In practice $W^+$ is computed from a basis of $S_{3/2}(\Gamma_0(M))$ to a finite precision $B$,
imposing the vanishing at the indices $n\le B$; this gives a space containing $W^+$, so the test
stays valid, and for $B$ large enough it is $W^+$ itself. The pairing with $V$ needs the
coefficients of $U(4)g$ up to the largest pole, i.e.\ $B/4$ at least that.

\section{Proof of Theorem~\ref{thm:main}}

\begin{proof}
The curve $C$ has genus three, and its Jacobian is isogenous to $A_f$ with $f=\texttt{609.2.a.d}$:
the characteristic polynomials of $T_5$, $T_{11}$, $T_{13}$ on $\Jac(C)$, read off the point counts
over $\F_p,\F_{p^2},\F_{p^3}$ from the trace formula, are
\[
x^3+6x^2+8x+2,\quad x^3+10x^2+28x+20,\quad x^3+6x^2+12x+8,
\]
the first irreducible, so $\Jac(C)$ is isotypic for a single newform orbit of dimension three; by
Jacquet--Langlands that orbit has level $21m$ with $m\mid58$ and is new at $3$ and $7$, and among the newform orbits of weight two and dimension three at the levels $21m$, $m\mid58$, that
Jacquet--Langlands allows, exactly one has these Hecke polynomials; it has level $609$ and
Atkin--Lehner signs $-,-,+$ at $3,7,29$, which become $+,+,+$ on the Shimura curve
(\texttt{tests/Hyperelliptic21\_58.m}). So Lemma~\ref{lem:inj} applies, with $X'^*=X_0^{21}(29)^*$.
The eight discriminants of the four fibres each carry one rational point of $C$ (the degree of the
field of definition of the CM points is one, and they form one Galois orbit). The torsion of
$\Jac(C)(\Q)$ has order dividing $\#\Jac(C)(\F_p)$ for every odd good prime, and the greatest common
divisor over the good primes up to $60$ is $2$.

The dual test of Theorem~\ref{thm:test} at $M=2436$, with the $233$ cusp forms computed to
precision $1600$, gives $\dim W^+=3$, the rank of $\Jac(C)(\Q)$, and certifies that
$F_2-F_4$, $F_2-F_5$, $F_2-F_6$ and $F_4-F_5$ are torsion, while refusing the differences
$P_{-16}-P_{-120}$, $P_{-88}-P_{-100}$, $P_{-16}-P_{-88}$, $P_{-112}-P_{-372}$, $P_{-16}-P_{-168}$
of points that are not fibres (\texttt{tests/\_offline/Hyperelliptic21\_58.m}). So the classes of
$F_2,F_4,F_5,F_6$ lie in one coset of a group of order at most $2$, two of them coincide, and $C$ has
two distinct effective divisors of degree two in one class: a $g^1_2$.
\end{proof}

\begin{remark}
In the pole-at-infinity dictionary \eqref{eq:div}, an odd discriminant $d$ and $4d$ are reached from
the same indices $n=|d|s^2$, and when $2\mid N$ both carry CM points, so $P_d$ and $P_{4d}$ cannot be
prescribed separately. This is why only the fibres in even discriminants enter the proof; the fifth
and sixth fibres of the numerical screen, $\{P_{-7},P_{-312}\}$ and $\{P_{-28},P_{-840}\}$, involve
$-7$ and $-28$ and are not used. The same obstruction leaves $X_0^{77}(10)^*$ with two usable fibres
against a torsion bound of $2$, and $X_0^{39}(34)^*$ with one against a bound of $4$.
\end{remark}

\begin{thebibliography}{9}
\bibitem{borcherds} R.~E.~Borcherds, \emph{The Gross--Kohnen--Zagier theorem in higher dimensions},
Duke Math. J. 97 (1999), 219–233; correction, Duke Math. J. 105 (2000), 183–184.
\bibitem{ez} M.~Eichler and D.~Zagier, \emph{The theory of Jacobi forms}, Progr. Math. 55,
Birkh\"auser, 1985.
\bibitem{gy} J.~Guo and Y.~Yang, \emph{Equations of hyperelliptic Shimura curves}, Compos. Math.
153 (2017), 1–40.
\bibitem{kohnen} W.~Kohnen, \emph{Newforms of half-integral weight}, J. reine angew. Math. 333
(1982), 32–72.
\bibitem{scheithauer} N.~R.~Scheithauer, \emph{The Weil representation of $\mathrm{SL}_2(\Z)$ and
some applications}, Int. Math. Res. Not. 2009, 1488–1545.
\end{thebibliography}
\end{document}
